\documentclass[11pt]{article}
\usepackage[margin=0.9in]{geometry}
\usepackage{amsmath,amssymb,amsthm}
\usepackage{enumitem}
\usepackage{booktabs}
\usepackage{array}
\usepackage[most]{tcolorbox}
\usepackage{hyperref}
\hypersetup{colorlinks=true,linkcolor=blue!50!black,urlcolor=blue!50!black,pdfpagemode=UseNone}

\newtheorem{theorem}{Theorem}
\setlength{\parskip}{4pt}
\setlength{\parindent}{0pt}

% Boxes used throughout: a yellow "picture it" box and a red "fix" box.
\newtcolorbox{intuition}{colback=yellow!12,colframe=yellow!55!black,
  title=\textbf{Picture it},fonttitle=\small,boxrule=0.6pt,left=6pt,right=6pt,top=3pt,bottom=3pt}
\newtcolorbox{fixbox}{colback=green!8,colframe=green!45!black,
  title=\textbf{Takeaway and fix},fonttitle=\small,boxrule=0.6pt,left=6pt,right=6pt,top=3pt,bottom=3pt}

\title{\textbf{Six RSA Attacks:\\ Håstad, Common Modulus, Bleichenbacher, Shared Primes,\\ Fermat, and Pollard's $p-1$}}
\author{}
\date{}

\begin{document}
\maketitle
\vspace{-2.5em}

Every attack below has the same shape: RSA is used with a small mistake in \emph{how it is deployed or how its keys are generated}, and a tool from Lesson 1 --- the Chinese Remainder Theorem, the Extended Euclidean Algorithm, RSA's multiplicative malleability, the plain gcd, the difference of two squares, or Fermat's little theorem --- turns that mistake into a full break. The first four attacks never factor $n$ at all. The last two do factor $n$, but only because a flawed key generator made $n$ much easier to factor than a product of two independent random primes. None of them breaks correctly generated and correctly used RSA.

\medskip
\renewcommand{\arraystretch}{1.25}
\begin{tabular}{@{}>{\raggedright\arraybackslash}p{2.6cm}>{\raggedright\arraybackslash}p{4.9cm}>{\raggedright\arraybackslash}p{3.2cm}>{\raggedright\arraybackslash}p{3.6cm}@{}}
\toprule
\textbf{Attack} & \textbf{The mistake} & \textbf{Tool used} & \textbf{Fix} \\
\midrule
1. Håstad & Same message, small $e$, many recipients & CRT & Randomized padding (OAEP) \\
2. Common modulus & One $n$ shared by several users & Extended Euclid & Never share $n$ \\
3. Bleichenbacher & Server reveals ``padding valid?'' & Malleability & OAEP, uniform errors \\
4. Shared prime & Bad randomness reuses one prime $p$ & gcd (Euclid) & Good entropy at key generation \\
5. Fermat & Primes $p,q$ chosen too close together & Difference of squares & Pick $p,q$ independently at random \\
6. Pollard $p-1$ & $p-1$ has only small prime factors & Fermat's little theorem + gcd & Primes with a large prime factor in $p-1$ \\
\bottomrule
\end{tabular}

\bigskip
\hrulefill

%=====================================================================
\section*{1. Håstad's Broadcast Attack}

\begin{intuition}
Think of $m^3 \bmod n$ as the reading on a clock with $n$ hours. If $m^3$ is bigger than $n$, the hand goes around the dial and you only see where it ended up --- that wraparound is what hides $m$. Now suppose the same $m^3$ is shown to you on \emph{three different clocks}, of sizes $n_1,n_2,n_3$. CRT lets you merge them into one giant clock of size $n_1n_2n_3$. If $m^3$ is smaller than that giant clock, the hand never wrapped, so the reading \emph{is} $m^3$, and a plain cube root gives $m$. Small $e$ means too little wraparound to hide anything.
\end{intuition}

\textbf{When this arises.} One message must go to several people: a vendor sends the same license key or firmware secret to many customers, or a group announcement is sent encrypted to each member. Every recipient has their own independent key pair, but many use the same tiny exponent (historically $e=3$, chosen because it makes encryption fast on weak hardware). The sender encrypts the identical $m$ separately under each recipient's public key.

\textbf{Setup.} Exponent $e=3$, the same message $m$ sent to three recipients with pairwise-coprime moduli $n_1,n_2,n_3$. The attacker sees
\[
c_i = m^3 \bmod n_i, \qquad i=1,2,3.
\]

\begin{theorem}
If $m<\min(n_1,n_2,n_3)$, an eavesdropper who sees $c_1,c_2,c_3$ recovers $m$ with no private key.
\end{theorem}

\textbf{Why it works.}
\begin{enumerate}[nosep]
\item CRT combines the three congruences $x\equiv c_i \pmod{n_i}$ into one unique answer $C$ modulo $N=n_1n_2n_3$. This uses only public data.
\item The true value $m^3$ satisfies all three congruences too, so by uniqueness $C\equiv m^3 \pmod N$.
\item Because $m<n_i$ for each $i$, we have $m^3<n_1n_2n_3=N$. So $C$ and $m^3$ are both in $[0,N)$ and congruent mod $N$, hence \emph{equal as ordinary integers}.
\item Take an ordinary integer cube root of $C$. Done.
\end{enumerate}

\textbf{Worked example.} Take $n_1=143=11\cdot13$, $n_2=323=17\cdot19$, $n_3=667=23\cdot29$, and secret message $m=17$. The three ciphertexts are
\[
c_1=17^3\bmod143=51,\qquad c_2=17^3\bmod323=68,\qquad c_3=17^3\bmod667=244.
\]
Each number alone looks unrelated to 17. The attacker now runs CRT with $N=143\cdot323\cdot667=30{,}808{,}063$. For each $i$, let $N_i=N/n_i$ and $y_i=N_i^{-1}\bmod n_i$:
\begin{center}
\begin{tabular}{@{}cccc@{}}
\toprule
$n_i$ & $c_i$ & $N_i$ & $y_i$ \\
\midrule
143 & 51 & 215{,}441 & 112 \\
323 & 68 & 95{,}381 & 286 \\
667 & 244 & 46{,}189 & 221 \\
\bottomrule
\end{tabular}
\end{center}
\[
C = 51\cdot215441\cdot112 + 68\cdot95381\cdot286 + 244\cdot46189\cdot221 \bmod N = 4913.
\]
Since $17^3=4913$ is far less than $N\approx 3\times10^7$, there was no wraparound at all, and $\sqrt[3]{4913}=17=m$.

\begin{fixbox}
Never send the same (or related) message to several recipients under a small $e$ without randomized padding. With OAEP each ciphertext encrypts a \emph{different} random-looking padded value, so the three equations no longer share one $m^3$. (In general, with exponent $e$ the attacker needs $e$ recipients.)
\end{fixbox}

\hrulefill

%=====================================================================
\section*{2. Common Modulus Attack}

\begin{intuition}
This is the classic water-jug puzzle. With a 17-litre jug and a 7-litre jug you can measure exactly 1 litre: fill the 7-litre jug five times and pour out the 17-litre jug twice ($5\cdot7-2\cdot17=1$). Here the ``jugs'' are the exponents. Ciphertext $c_1$ is $m$ raised to the power 17, and $c_2$ is $m$ raised to the power 7. Multiplying five copies of $c_2$ and ``removing'' two copies of $c_1$ leaves $m$ raised to the power $5\cdot 7-2\cdot 17=1$, that is $m$ itself.
\end{intuition}

\textbf{When this arises.} Generating a fresh modulus $n=pq$ for every user is slow, because finding large primes dominates the cost. A careless authority might generate $n$ once and hand each user a different exponent pair $(e_i,d_i)$ over that same $n$, thinking the different exponents keep users separate. This was a real anti-pattern in some early multi-user and embedded key-management designs.

\textbf{Setup.} One shared modulus $n$; two users with $\gcd(e_1,e_2)=1$; the same message sent to both:
\[
c_1=m^{e_1}\bmod n,\qquad c_2=m^{e_2}\bmod n.
\]

\begin{theorem}
If $\gcd(e_1,e_2)=1$, anyone who knows $(n,e_1,e_2,c_1,c_2)$ recovers $m$ with no private key.
\end{theorem}

\textbf{Why it works.}
\begin{enumerate}[nosep]
\item Extended Euclid on $(e_1,e_2)$ gives integers $a,b$ with $ae_1+be_2=1$.
\item Then $c_1^{\,a}c_2^{\,b}\equiv m^{ae_1}m^{be_2}=m^{ae_1+be_2}=m \pmod n$, by exponent laws alone.
\item One of $a,b$ is negative. For a negative exponent, first compute the modular inverse of the corresponding ciphertext (again by Extended Euclid), then raise to the absolute value.
\end{enumerate}
Notice that the attacker never factors $n$ and never touches $d_1$ or $d_2$.

\textbf{Worked example.} The classic toy key $n=3233=61\cdot53$, with $e_1=17$, $e_2=7$ and secret $m=42$:
\[
c_1=42^{17}\bmod3233=2557,\qquad c_2=42^{7}\bmod3233=240.
\]
\emph{Step 1 --- Euclid on $(17,7)$.}
$17=2\cdot7+3$, $\ 7=2\cdot3+1$. Back-substituting: $1=7-2\cdot3=7-2(17-2\cdot7)=5\cdot7-2\cdot17$. So $a=-2$, $b=5$.

\emph{Step 2 --- handle the negative exponent.} $c_1^{-1}\bmod3233=3123$ (check: $2557\cdot3123\equiv1$).

\emph{Step 3 --- combine.}
\[
m\equiv (c_1^{-1})^{2}\cdot c_2^{5} \equiv 2401\cdot1904 \equiv 42 \pmod{3233}.
\]

\textbf{An even worse consequence.} Any \emph{insider} who legitimately owns one pair $(e_1,d_1)$ on the shared $n$ can use the fact that $e_1d_1-1$ is a multiple of $\varphi(n)$ to factor $n$. After that, every other user's private key is compromised too, not just one message.

\begin{fixbox}
Never share a modulus across users or purposes, even with different exponents. It is broken as soon as two exponents on it are coprime, which is the typical case.
\end{fixbox}

\hrulefill

%=====================================================================
\section*{3. Bleichenbacher's Padding Oracle Attack}

\begin{intuition}
Picture the hidden number $m$ as a runner on a circular track of length $n$. A bouncer stands at one narrow gate (the window $[2B,3B)$) and answers only ``yes, someone is in the gate'' or ``no''. You cannot see $m$. But you can make the runner go $s$ times faster than $m$ does --- that is just multiplying by $s$ --- and ask whether the sped-up runner is now in the gate. A ``yes'' means the true $m$ lies in one of a few thin, evenly spaced slices of the track. Choose the next $s$ so that only one slice remains consistent, repeat, and the slices shrink like a game of Twenty Questions until only $m$ is left.
\end{intuition}

\textbf{When this arises.} A server decrypts RSA ciphertexts that were padded before encryption, and --- by an error message, a dropped connection, or even a timing difference --- lets the sender tell whether the padding was valid. Nobody intends to leak anything else, but it is very easy to do by accident: TLS servers using the PKCS\#1 v1.5 padding standard commonly sent a distinguishable ``bad padding'' alert. This is not history. Refinements (Bardou et al., 2012; the ROBOT attack, 2017) found major production web servers still exploitable nearly twenty years after the original 1998 paper, and it remains one of the most consequential vulnerabilities ever found in a widely deployed protocol.

\textbf{What the oracle checks.} PKCS\#1 v1.5 formats a $k$-byte block as
\[
\text{EB} = \texttt{0x00}\ \|\ \texttt{0x02}\ \|\ \text{PS}\ \|\ \texttt{0x00}\ \|\ D,
\]
where $D$ is the message and PS is nonzero random padding filling the block to $k$ bytes. For example, with $k=8$, $D=\texttt{0x1234}$ and $\text{PS}=\texttt{0xA1\,7F\,3C}$:
\[
\text{EB}=\texttt{00\ 02\ A1\ 7F\ 3C\ 00\ 12\ 34}.
\]
A \textbf{padding oracle} answers one yes/no question about a ciphertext $c'$ of the attacker's choosing: \emph{does $(c')^d \bmod n$ start with \texttt{00 02}?}

Read as a big integer, ``starts with \texttt{00 02}'' simply means the value lies in $[2B,3B)$, where $B=256^{k-2}$. That is a very narrow window: a random ciphertext passes only about 1 time in $65{,}536$. The attack is a way of asking this same narrow question about \emph{cleverly chosen} values instead of random ones.

\textbf{The two ingredients.}
\begin{itemize}[nosep]
\item \emph{Malleability.} From $c=m^e\bmod n$ anyone can compute $c\cdot s^e\bmod n$ for any chosen $s$. This is a valid encryption of $m\cdot s \bmod n$, and it needs no secret.
\item \emph{A YES is a strong clue.} If the oracle says YES for $c\cdot s^e$, then $ms \bmod n\in[2B,3B)$ for an $s$ the attacker knows.
\end{itemize}

\textbf{From one YES to a small set of intervals for $m$.} Unwrapping ``mod $n$'', a YES for $s$ means $ms = rn + t$ with $t \in [2B,3B)$ for some unknown whole number $r$ (how many laps the fast runner made). Dividing by $s$:
\[
m\in\left[\frac{2B+rn}{s},\ \frac{3B-1+rn}{s}\right]\qquad\text{for some integer } r.
\]
Early on several values of $r$ are all consistent, so one YES can leave more than one candidate interval (this really happens, as the example shows). Later YES/NO answers eliminate the wrong ones.

\textbf{The algorithm in three moves.}
\begin{enumerate}[nosep]
\item \textbf{Get started.} If $c$ is not itself oracle-valid, find a small $s_0$ with $c\cdot s_0^e$ oracle-valid and work with $ms_0$ instead. (If $c$ is already validly padded, $s_0=1$.)
\item \textbf{Find the next $s_i$.} While several intervals survive, try $s$ upward one at a time until the oracle says YES. Once only one interval $[a,b]$ survives, jump straight to $s$ values computed from its endpoints; this roughly doubles $s$ each round and halves the interval.
\item \textbf{Shrink.} For each surviving interval and each possible lap count $r$, compute the new sub-interval from the formula above, keep the non-empty ones, and merge overlaps. Stop when one interval containing a single integer remains. That integer is EB (undo $s_0$ if it was used).
\end{enumerate}

\textbf{Worked example (computed and verified).} Toy parameters: $n=16{,}826{,}323$ ($p=4093$, $q=4111$), $e=17$, $d=10{,}882{,}313$, block length $k=4$, so $B=256^2=65{,}536$ and the valid window is $[131{,}072,\ 196{,}608)$. The secret is $m=150{,}000$ (hex \texttt{0x0249F0}, which starts \texttt{00 02}, so $s_0=1$), with ciphertext $c=8{,}054{,}451$.

\emph{The first query in detail.} The attacker tries $s=1,2,3,\dots$ and the first YES comes at $s_1=450$. Why? Because
\[
450\cdot150{,}000 = 67{,}500{,}000 = 4\cdot n + 194{,}708,
\]
and $194{,}708$ lies in the window. The attacker does not know this, but sees the YES and reasons: ``$450m$ made $r$ laps and landed in the window.'' Two lap counts fit the information so far, $r=4$ and $r=5$, giving the two intervals in row 1. The true one is $r=4$.

\begin{center}
\begin{tabular}{@{}rrll@{}}
\toprule
$i$ & $s_i$ & surviving interval(s) for $m$ & width \\
\midrule
0 & --- & $[131{,}072,\ 196{,}607]$ & (initial) \\
1 & 450 & $[149{,}859,\,150{,}004]$ and $[187{,}251,\,187{,}396]$ & two intervals \\
2 & 562 & $[149{,}934,\ 150{,}004]$ & 70 \\
3 & 1{,}235 & $[149{,}977,\ 150{,}004]$ & 27 \\
4 & 2{,}581 & $[149{,}995,\ 150{,}004]$ & 9 \\
5 & 5{,}610 & $[149{,}995,\ 150{,}002]$ & 7 \\
6 & 11{,}331 & $[149{,}995,\ 150{,}000]$ & 5 \\
7 & 22{,}885 & $[149{,}998,\ 150{,}000]$ & 2 \\
8 & 45{,}881 & $[149{,}999,\ 150{,}000]$ & 1 \\
9 & 91{,}873 & $[150{,}000,\ 150{,}000]$ & \textbf{0 --- recovered} \\
\bottomrule
\end{tabular}
\end{center}

Two things to point out to students. First, the false interval from $r=5$ is killed at $i=2$: the next YES ($s=562$) is consistent only with the true interval. Second, once one interval remains, $s_i$ roughly doubles each round ($562\to1235\to2581\to5610\to\cdots$) and the interval width roughly halves: this is binary search. Nine oracle queries against a modulus above sixteen million pinned $m$ down exactly.

\textbf{A caveat about the toy numbers.} Here $n$ is a 24-bit number, so the block's leading \texttt{00} byte is automatic and the oracle really checks only the second byte (about 1 random value in 256 passes). With a real 2048-bit modulus both bytes must match (about 1 in 65{,}000), so a real attack needs thousands to about a million queries rather than nine. The mechanics are identical.

\begin{fixbox}
Never let the reason for a decryption failure be distinguishable --- by error content, timing, or connection behavior. The complete fix is to stop using PKCS\#1 v1.5 encryption in new designs and use OAEP (Lesson 3), whose security proof rules out this whole class of attack.
\end{fixbox}

\hrulefill

%=====================================================================
\section*{4. Shared-Prime (Batch GCD) Attack}

\begin{intuition}
Factoring a single 2048-bit number is like trying to un-bake a cake: hopeless. But if two cakes were baked with a common ingredient, you can find it just by comparing them --- and Euclid's algorithm does that comparison in a fraction of a millisecond. If two RSA moduli happen to share a prime, computing $\gcd(n_1,n_2)$ hands you that prime, and dividing gives the other prime of each. Both keys are broken, and neither owner did anything obviously wrong.
\end{intuition}

\textbf{When this arises.} Devices such as routers, firewalls and other embedded systems often generate their RSA key the first time they boot. At that moment the operating system's random number generator has had almost no entropy (no disk, no keyboard, no network jitter yet), so many devices produce \emph{the same} first prime $p$. A little entropy trickles in (a clock tick, a packet) before the second prime is chosen, so their $q$'s differ. The result is many public keys of the form $n_i=p\,q_i$ with a common $p$.

This is not hypothetical. In 2012 two independent teams (Lenstra et al., \emph{``Ron was wrong, Whit is right''}; Heninger et al., \emph{``Mining your Ps and Qs''}) collected millions of public RSA keys from the internet and factored thousands of them this way, with no ciphertext and no private key needed --- just public moduli. To avoid comparing every pair, they used ``batch GCD'', a product-tree algorithm that checks millions of moduli against each other in hours.

\textbf{Setup.} The attacker has two public moduli
\[
n_1=p\,q_1,\qquad n_2=p\,q_2,\qquad q_1\ne q_2,
\]
and knows nothing else about the factorizations.

\begin{theorem}
Then $\gcd(n_1,n_2)=p$. Consequently $q_1=n_1/p$ and $q_2=n_2/p$, and the private keys of \emph{both} moduli can be computed.
\end{theorem}

\textbf{Why it works.}
\begin{enumerate}[nosep]
\item $p$ divides both $n_1$ and $n_2$, so $p\mid\gcd(n_1,n_2)$.
\item Since $q_1$ and $q_2$ are distinct primes (and differ from $p$), $\gcd(n_1,n_2)=p\cdot\gcd(q_1,q_2)=p$.
\item With all primes known, compute $\varphi(n_1)=(p-1)(q_1-1)$ and $d_1=e_1^{-1}\bmod\varphi(n_1)$ --- exactly the recipe the legitimate key owner used. The same works for $n_2$.
\end{enumerate}
The gcd takes $O(\log^2 n)$ bit operations, whereas factoring a 2048-bit modulus directly is infeasible. That gap is the entire attack.

\textbf{Worked example.} Alice's key is $n_1=3233=61\cdot53$ and Bob's key is $n_2=3599=61\cdot59$; both use $e=17$. Eve sees only the two public moduli, so all she starts with is $3233$ and $3599$.

\emph{Step 1 --- Euclid.}
\begin{align*}
3599 &= 1\cdot3233+366,\\
3233 &= 8\cdot366+305,\\
366 &= 1\cdot305+61,\\
305 &= 5\cdot61+0,
\end{align*}
so $\gcd(3233,3599)=61$. Eve has found the shared prime $p=61$ in four division steps.

\emph{Step 2 --- factor both keys.} $q_1=3233/61=53$ and $q_2=3599/61=59$.

\emph{Step 3 --- rebuild Alice's private key.} $\varphi(n_1)=60\cdot52=3120$ and $d_1=17^{-1}\bmod3120=2753$ (check: $17\cdot2753=46801=15\cdot3120+1$). Bob's key falls the same way: $\varphi(n_2)=60\cdot58=3480$, $d_2=1433$.

\emph{Step 4 --- read a message.} Suppose someone sent Alice the ciphertext $c=2183$. Eve computes $2183^{2753}\bmod3233=1234$, the plaintext, without ever being given a key.

\begin{fixbox}
Generate keys only after the system has gathered enough entropy, use the operating system's cryptographic random number generator, and draw $p$ and $q$ from fresh randomness. Avoid generating keys on the first boot of a device with no entropy source. Operators can also check their own newly generated keys against large public key collections for shared factors. Notice the lesson: RSA's security depends not only on the math but on the quality of the random numbers used to pick the primes.
\end{fixbox}

\hrulefill

%=====================================================================
\section*{5. Fermat's Factorization (Close Primes)}

\begin{intuition}
Think of $n=pq$ as the area of a $p\times q$ rectangle of unit tiles. Any rectangle can be rearranged into a big square with a smaller square cut out of one corner. Let $a$ be the \emph{midpoint} of $p$ and $q$ and $b$ the \emph{distance from the midpoint to each}, so $p=a-b$ and $q=a+b$. Then the rectangle has the same area as a square of side $a$ minus a square of side $b$:
\[ pq=(a-b)(a+b)=a^2-b^2. \]
If $p$ and $q$ are nearly equal, the rectangle is almost a square already, so $a$ is only slightly above $\sqrt n$ and the cut-out square is tiny. The attacker simply tries $a=\lceil\sqrt n\rceil,\ \lceil\sqrt n\rceil+1,\dots$ until $a^2-n$ turns out to be a perfect square. That is a search from a known starting point, not a blind one.
\end{intuition}

\textbf{When this arises.} A key generator picks one random prime $p$ and then, to save time or because of a flawed random source, picks $q$ close to it --- for instance the next prime after $p$, or $p$ plus a small random offset. The idea itself dates to 1643, and for a long time it was regarded as a textbook curiosity. In 2022 Hanno Böck applied it to large collections of real public keys and found vulnerable ones: some Canon and Fujifilm (Fuji Xerox) printers generated keys with a cryptographic module that did not randomize its two primes enough (the Canon issue is tracked as CVE-2022-26351). It is a rare flaw, but it is real, and it needs nothing except the public modulus.

\textbf{Setup.} The attacker knows only $n=pq$ with $p<q$ odd primes. Nothing else.

\begin{theorem}
Let $a=\tfrac{p+q}{2}$ and $b=\tfrac{q-p}{2}$ (integers, since $p,q$ are odd). Then $a^2-n=b^2$, and $a$ is the smallest integer $\ge\sqrt n$ for which $a^2-n$ is a perfect square. Testing $a=\lceil\sqrt n\rceil,\lceil\sqrt n\rceil+1,\dots$ therefore finds $a$ after at most
\[
\frac{(q-p)^2}{8\sqrt n}+1
\]
tries, and then $p=a-b$, $q=a+b$.
\end{theorem}

\textbf{Why it works.}
\begin{enumerate}[nosep]
\item Algebra: $a^2-b^2=(a-b)(a+b)=p\cdot q=n$, so $a^2-n=b^2$ is a perfect square.
\item Conversely, if $a^2-n=b^2$ for any integers $a>b\ge0$, then $n=(a-b)(a+b)$ is a factorization. Since $n=pq$ has only one pair of nontrivial factors, the first hit gives exactly $p$ and $q$.
\item How many tries? The search starts at $\lceil\sqrt n\rceil\ge\sqrt n$ and stops at $a=\tfrac{p+q}{2}$, so it takes at most $\tfrac{p+q}{2}-\sqrt{pq}+1$ steps. Now
\[
\frac{p+q}{2}-\sqrt{pq}=\frac{(\sqrt q-\sqrt p)^2}{2}=\frac{(q-p)^2}{2(\sqrt q+\sqrt p)^2}\le\frac{(q-p)^2}{8\sqrt n},
\]
because $(\sqrt p+\sqrt q)^2=p+q+2\sqrt n\ge4\sqrt n$ (as $p+q\ge2\sqrt{pq}$).
\end{enumerate}
So the work depends on the \emph{gap} $q-p$, and grows with its square.

\textbf{Worked example.} Eve sees the public key $n=5959$, $e=17$. Here $\sqrt{5959}\approx77.2$, so she starts at $a=78$:
\begin{center}
\begin{tabular}{@{}rrrl@{}}
\toprule
$a$ & $a^2$ & $a^2-5959$ & perfect square? \\
\midrule
78 & 6084 & 125 & no ($11^2=121,\ 12^2=144$) \\
79 & 6241 & 282 & no ($16^2=256,\ 17^2=289$) \\
80 & 6400 & 441 & \textbf{yes}, $441=21^2$ \\
\bottomrule
\end{tabular}
\end{center}
So $a=80$, $b=21$, and $p=80-21=59$, $q=80+21=101$. Check: $59\cdot101=5959$. Three tries, and no factoring in the usual sense.

\emph{Finishing the break.} $\varphi(n)=58\cdot100=5800$ and $d=17^{-1}\bmod5800=3753$ (check: $17\cdot3753=63801=11\cdot5800+1$). If Alice's ciphertext is $c=2029$, Eve computes $2029^{3753}\bmod5959=1234$, the plaintext.

\emph{Contrast: primes that are far apart.} Take $p=101$, $q=1009$, so $n=101909$ and $\sqrt n\approx319.2$. Now $a$ must climb all the way to $\tfrac{101+1009}{2}=555$ (with $b=454$), which is $236$ tries, against a bound of $\tfrac{908^2}{8\sqrt n}+1\approx323$. The bigger the gap, the longer the climb.

\textbf{How big a gap is safe?} For a 2048-bit modulus, $8\sqrt n\approx2^{1027}$. The bound above then says:
\begin{center}
\begin{tabular}{@{}>{\raggedright\arraybackslash}p{6.8cm}>{\raggedright\arraybackslash}p{7.6cm}@{}}
\toprule
\textbf{Gap $q-p$} & \textbf{Tries needed (at most)} \\
\midrule
below about $2^{512}$ & 1 \\
$2^{512+t}$ & about $2^{2t}/8$ (each extra bit of gap quadruples the work) \\
about $2^{1023}$ (two independent random 1024-bit primes) & about $2^{1019}$, hopeless \\
\bottomrule
\end{tabular}
\end{center}
So Fermat only works when the two primes agree in roughly their top half of bits. Independent random primes never do.

\textbf{Try it (about 10 lines of Python).}
\begin{verbatim}
from math import isqrt

def fermat(n):
    a = isqrt(n)
    if a * a < n:            # a = ceil(sqrt(n))
        a += 1
    while True:
        b2 = a * a - n
        b = isqrt(b2)
        if b * b == b2:      # a^2 - n is a perfect square
            return a - b, a + b
        a += 1
\end{verbatim}
\texttt{fermat(5959)} returns \texttt{(59, 101)}. A good classroom demo: generate two $512$-bit primes with \texttt{p = nextprime(r)} and \texttt{q = nextprime(p + 10**9)}. The $1024$-bit modulus falls on the very first try.

\begin{fixbox}
Choose $p$ and $q$ \emph{independently} from a good random source, and reject the pair if they are too close. Standards such as FIPS~186-4 require $|p-q|>2^{\,n_{\text{len}}/2-100}$. With independent random primes this is satisfied automatically, so the attack only ever hits a faulty generator, never a correct one.
\end{fixbox}

\hrulefill

%=====================================================================
\section*{6. Pollard's $p-1$ Factoring Method}

\begin{intuition}
Look at the powers of a number $a$ on a clock with $p$ hours. By Fermat's little theorem the sequence $a,a^2,a^3,\dots$ lands back on $1$ exactly when the exponent is a multiple of (the period dividing) $p-1$. Now suppose $p-1$ is built only from small steps, such as $2^4\cdot3^2\cdot7$. Then a number like $B!=1\cdot2\cdot3\cdots B$ is a multiple of $p-1$, so $a^{B!}$ reads $1$ on the $p$-clock. The attacker cannot see the $p$-clock --- all she has is $n=pq$ --- but she can compute $a^{B!}\bmod n$ anyway. The other prime $q$ has its own, much longer period $q-1$, so on the $q$-clock the same number almost surely does \emph{not} read $1$. Subtracting $1$ and taking a gcd with $n$ cleanly separates the two clocks and reveals $p$.
\end{intuition}

\textbf{Vocabulary.} A number is \emph{$B$-smooth} if all its prime factors are at most $B$. For example $1008=2^4\cdot3^2\cdot7$ is $7$-smooth, while $1018=2\cdot509$ is not even $100$-smooth.

\textbf{When this arises.} A key generator builds a prime from small pieces, say $p=2m+1$ where $m$ is a product of small primes, because such a form makes it easy to prove $p$ prime. The resulting $p-1$ is smooth. Standards from the 1990s therefore asked for ``strong primes'', and modern ones such as FIPS~186 still require $p-1$ to contain a large prime factor. John Pollard published the method in 1974, and it remains one of the first things factoring software tries on a number, because it costs almost nothing.

\textbf{Setup.} The attacker knows only $n=pq$ and picks a bound $B$ (a guess about how smooth $p-1$ might be).

\begin{theorem}
Let $M=B!$ and $a$ coprime to $n$. Suppose $p-1$ divides $M$ but $a^M\not\equiv1\pmod q$. Then
\[
\gcd\bigl(a^M-1,\ n\bigr)=p.
\]
\end{theorem}

\textbf{Why it works.}
\begin{enumerate}[nosep]
\item Write $M=(p-1)t$. By Fermat's little theorem $a^{p-1}\equiv1\pmod p$, so $a^M=(a^{p-1})^t\equiv1\pmod p$. Hence $p\mid a^M-1$, and $p\mid n$, so $p\mid\gcd(a^M-1,n)$.
\item By assumption $a^M\not\equiv1\pmod q$, so $q\nmid a^M-1$. The gcd is therefore a divisor of $n=pq$ that contains $p$ but not $q$, so it equals $p$.
\item The attacker never needs $p$ to compute this: $a^M\bmod n$ is an ordinary modular exponentiation, and $\gcd\bigl((a^M\bmod n)-1,\,n\bigr)$ equals $\gcd(a^M-1,n)$.
\end{enumerate}

\textbf{The algorithm.} Start with $a=2$. For $j=2,3,4,\dots$ replace $a$ by $a^j\bmod n$. After step $j$ we have $a\equiv2^{j!}\pmod n$, because $\bigl(2^{(j-1)!}\bigr)^j=2^{j!}$. After each step compute $d=\gcd(a-1,n)$:
\begin{itemize}[nosep]
\item $d=1$: not yet; keep going (up to $B$).
\item $1<d<n$: success, $d$ is a prime factor.
\item $d=n$: both $p-1$ and $q-1$ divide $j!$ at the same time. Retry with a different base $a$.
\end{itemize}

\textbf{Worked example 1 (by hand).} $n=299$. In fact $299=13\cdot23$, with $13-1=12=2^2\cdot3$ (smooth) and $23-1=22=2\cdot11$. Eve does not know this.
\begin{center}
\begin{tabular}{@{}cccc@{}}
\toprule
$j$ & $a=2^{j!}\bmod299$ & $a-1$ & $\gcd(a-1,299)$ \\
\midrule
2 & $2^2=4$ & 3 & 1 \\
3 & $4^3=64$ & 63 & 1 \\
4 & $64^4\equiv27$ & 26 & \textbf{13} \\
\bottomrule
\end{tabular}
\end{center}
(For $j=4$: $64^2=4096\equiv209$, and $209^2=43681\equiv27\pmod{299}$.) At $j=4$ we have $j!=24$, and $12\mid24$. So $2^{24}\equiv1\pmod{13}$; indeed $27=2\cdot13+1$. But $22\nmid24$, and indeed $27\equiv4\pmod{23}$, not $1$. The gcd isolates $13$, and $299/13=23$.

\textbf{Worked example 2 (a small RSA key).} $n=1028171=1009\cdot1019$, $e=17$. The hidden structure is $1009-1=1008=2^4\cdot3^2\cdot7$ and $1019-1=1018=2\cdot509$. Since $1008\mid7!=5040$ (indeed $5040=5\cdot1008$), the attack should succeed at $j=7$, but $509$ is a prime factor of $q-1$ that $7!$ cannot cancel.
\begin{center}
\begin{tabular}{@{}cccc@{}}
\toprule
$j$ & $a=2^{j!}\bmod n$ & $\gcd(a-1,n)$ & \\
\midrule
2 & 4 & 1 & \\
3 & 64 & 1 & \\
4 & 326\,480 & 1 & \\
5 & 321\,438 & 1 & \\
6 & 709\,177 & 1 & ($720$ is not a multiple of $1008$) \\
7 & 615\,491 & \textbf{1009} & ($5040=5\cdot1008$) \\
\bottomrule
\end{tabular}
\end{center}
So $p=1009$, $q=n/p=1019$, and with $\varphi(n)=1008\cdot1018$ Eve computes $d=e^{-1}\bmod\varphi(n)$ and owns the key. Seven modular exponentiations did what trial division by about $1000$ numbers would have done here, but the point is how it scales.

\textbf{Scaling up.} The work is only about $B$ small modular exponentiations, no matter how large $p$ is. In a test with a $754$-bit modulus, whose smaller prime $p$ ($242$ bits) had $p-1$ made of primes below $300$ and whose other prime was a random $512$-bit prime, the method returned $p$ at $j=281$ in a fraction of a second. By contrast, on a modulus with a $512$-bit \emph{safe} prime ($p=2p'+1$ with $p'$ prime), the same code with $B=3000$ found nothing.

\textbf{Try it (Python).}
\begin{verbatim}
from math import gcd

def pollard_p_minus_1(n, B=100000):
    a = 2
    for j in range(2, B + 1):
        a = pow(a, j, n)         # now a = 2^(j!) mod n
        d = gcd(a - 1, n)
        if 1 < d < n:
            return d             # a nontrivial factor
        if d == n:
            return None          # both p-1, q-1 smooth: change base
    return None                  # B too small
\end{verbatim}
\texttt{pollard\_p\_minus\_1(299)} returns \texttt{13}, and \texttt{pollard\_p\_minus\_1(1028171)} returns \texttt{1009}.

\textbf{What stops it.} If $p-1$ has a large prime factor $r$, the attacker needs $B\ge r$ before $B!$ becomes a multiple of $p-1$. For a safe prime $p=2p'+1$ this means $B\ge p'\approx p/2$, which costs as much as trial division and is hopeless. A random $1024$-bit prime has $p-1$ smooth up to $B=2^{20}$ with a probability far below $2^{-100}$, so smoothness is essentially never an accident. The danger comes only from generators that build $p$ so that $p-1$ is smooth, on purpose or by structure.

\begin{fixbox}
Generate primes so that $p-1$ (and $q-1$) contains a large prime factor, as FIPS~186 key generation does. Safe primes $p=2p'+1$ achieve this but are not required; random large primes already work with overwhelming probability. Together with Section~5 the lesson is the same: the key generator must produce $p$ and $q$ that are \emph{large, independent, and unstructured}, because RSA's security rests on $n$ looking like a random product of two large primes, not merely on $n$ being large.
\end{fixbox}

\hrulefill

%=====================================================================
\medskip
\textbf{Where each tool came from.} All six attacks reuse Lesson 1 material as \emph{attack} tools rather than construction tools: CRT (attack 1), the Extended Euclidean Algorithm (attack 2), multiplicative malleability (attack 3), the plain Euclidean gcd (attack 4), the difference of two squares (attack 5), and Fermat's little theorem together with the gcd (attack 6). Attacks 1 and 2 exploit a \emph{misuse} of correctly generated keys; attack 3 exploits an \emph{information leak}; attacks 4, 5 and 6 exploit \emph{flawed key generation}. Attacks 5 and 6 are the only ones that actually factor $n$, and they succeed only when $p$ and $q$ are not independent, unstructured, random primes.

\end{document}
